\section{Spectral and torsional comparison}\label{sec:consequences}

We now derive the model-ball variational comparisons stated in
Corollary~\ref{cor:model-ball-spectral}. Theorem~\ref{thm:main}
allows the modulus of each compactly supported smooth test on the
original domain to be replaced by an equimeasurable radial test on its
model ball with no greater energy.
The point to check is that the radial rearrangement belongs to the
model ball's Dirichlet Sobolev space.

\begin{proof}[Proof of Corollary~\ref{cor:model-ball-spectral}]
Let $M$, $D$, $\kappa$ and $p$ satisfy the hypotheses of the corollary.
Fix $0\ne v\in C_c^\infty(D)$ and extend $u=|v|$ by zero to $M$;
this is a compactly supported Lipschitz function.
Put $b=\sqrt{-\kappa}$ and take the interval density
$a(r)=\mathcal A_b(r)$ on $(0,\infty)$, whose cumulative mass is
$F_a(r)=\mathcal V_b(r)$. Thus $a(F_a^{-1}(s))=\mathcal I_b(s)$.
Apply Nobili and Violo's P\'olya--Szeg\H{o} theorem
\cite[Theorem~1.1]{NobiliViolo2025} on $M$ with constant one.
Its metric BV perimeter is the ambient Riemannian perimeter
\eqref{eq:ambient-perimeter}, as recalled after Definition~2.8 of
\cite{AntonelliBrueFogagnoloPozzetta2022}. Hence Theorem~\ref{thm:main}
gives the required support-neighborhood bound at every finite volume,
trivially also for sets of infinite perimeter. The positive superlevels
of $u$ have finite volume, and its relaxed metric $p$-energy is at most its Riemannian
energy \cite[Definition~2.1]{NobiliViolo2025}. The theorem and
\cite[Lemma~3.1]{NobiliViolo2025} give the decreasing equimeasurable
rearrangement $u^*$ for $a(r)\,dr$, bounded and locally absolutely
continuous, with weighted derivative energy at most
$\int_D|\nabla_gv|_g^p\,d\vol_g$.

Write $B=B_\kappa(|D|)$, with center $o$. Since $\supp u\Subset D$, its volume is strictly
below $|D|$, so $u^*$ vanishes on an interval before
$R=\mathcal V_b^{-1}(|D|)$.
We use the radial construction in
\cite[proof of Theorem~1.6, Section~7.1]{NobiliViolo2025}.
The polar density is $a(r)=n\omega_n r^{n-1}w_b(r)$, where
$w_b(r)=(\sn_b(r)/r)^{n-1}$ extends smoothly and positively at zero
and is log-convex. Define
$\widetilde u(x)=u^*(d_\kappa(o,x))$. In normal coordinates it is bounded
and locally $W^{1,p}$ away from the origin, with finite gradient
$p$-energy. Its gradient is therefore also locally integrable across
the origin. Integrating by parts outside the radius-$\varepsilon$ ball,
the inner boundary term is $O(\varepsilon^{n-1})$ because $u^*$ is
bounded. This term tends to zero since $n\ge2$, so the weak gradient
extends across the origin. Equivalence of Sobolev norms in smooth
coordinates and vanishing near $R$ now give
$\widetilde u\in W^{1,p}_0(B)$.
The model radial gradient is the radial derivative, and therefore
\[
 \int_B|\nabla_\kappa\widetilde u|_\kappa^p\,d\vol_\kappa
 \le\int_D|\nabla_gv|_g^p\,d\vol_g,
 \qquad
 \int_B|\widetilde u|^q\,d\vol_\kappa
 =\int_D|v|^q\,d\vol_g\quad(q=1,p).
\]
Taking the Rayleigh infimum and the torsion-quotient supremum proves the
claims. The defining density of $C_c^\infty(D)$ and the finite volume of
$D$ allow both variational values to be computed on the tests used above.
\end{proof}
